%% notes-tex/database/ecoli-datasets/sections/benchmark.tex
%% [[experiments.database.ecoli-datasets]]
%% https://github.com/Mjvolk3/torchcell/tree/main/notes-tex/database/ecoli-datasets/sections/benchmark.tex
%% SOURCE: hand-authored prose; every number quoted here is a macro from
%% tables/numbers.tex, generated by
%% experiments/database/scripts/build_ecoli_genslm_benchmark_table.py, and the pick
%% list is the PICKS registry in that script

\section{The short list for a gene-sequence encoding benchmark\secstatus{todo}}\label{sec:benchmark}

\subsection{What is being benchmarked\secstatus{todo}}

GenSLM (Zvyagin et al.\ 2023, mirror key \texttt{zvyaginGenSLMsGenomescaleLanguage2023})
tokenizes a gene as codons and was pretrained on more than 110 million prokaryotic gene
sequences from BV-BRC. \file{torchcell/datasets/genslm.py} embeds each gene's spliced
coding sequence, in frame from the start codon and truncated to the first 2,048 codons
when longer, and stores the masked mean of the last hidden layer: 512 values per gene
for the 25M checkpoint and 1,840 for the 250M one, both of which are on GilaHyper with
a sha256 manifest. The stores built so far are for \org{P. putida} KT2440; the
\org{E. coli} K-12 MG1655 and BW25113 genomes expose a coding sequence through the same
protocol, so their stores are a build away, not a code change.

The question a benchmark asks of such an encoding is narrow. Take a dataset whose label
attaches to a gene, hold out a set of genes, fit a small head on the vectors of the
remaining genes, and score the held-out ones. That score has meaning only against two
controls at the same width: a one-hot or random vector per gene, which carries gene
identity and nothing else, and a codon-frequency vector, which carries the one feature
of a coding sequence a language model is sure to have learned. A codon model earns its
place when it beats both; it is doing something different from either when it beats
one and not the other.

\subsection{The selection rule\secstatus{todo}}

Three conditions, applied in order to every row of Table~\ref{tab:ecoli-datasets}.

\pfstep{The label sits on one gene or one gene pair} The sequence-only fit tier must be
\emph{best}, or \emph{good} only because the gene is crossed with a condition, since the
condition is a second input the gene encoding does not have to carry. This removes
every whole-strain readout: the titers, the growth rates of an engineered strain, the
multi-omics panels of a wild type where the label is a profile of a strain rather than
of a gene. It also removes the \guideStoreCount{} stores whose label sits on a guide
RNA rather than on the gene it targets (Sec.~\ref{sec:caveats}).

\pfstep{Enough labeled genes for a held-out split} A benchmark scored on a few dozen
held-out genes cannot separate two encodings. Every core pick carries between
\coreGenesMin{} and \coreGenesMax{} genes with a label, read from a full scan of the
store, so a ten percent test fold holds hundreds of genes. The count that matters is
genes, not instances: a proteome of \schmidtEnv{} conditions is \schmidtInstances{}
instances and \schmidtGenes{} labeled genes, and it is the latter that is split.

\pfstep{One primary per task family} A short list is one dataset per kind of question,
with the datasets that pose the same question named beside it rather than listed on
their own. The families are a binary call per gene, a quantitative fitness per gene
and condition, a molecular level of the gene's own product, and a pairwise label.

\subsection{The picks\secstatus{todo}}

Table~\ref{tab:shortlist} is the list. The four core rows cover the four families and
should be run first; the four second-tier rows add a translation-level, a
transcript-level, an imaging and a knockdown task once the core four have been scored.

\pfstep{Essentiality, Goodall 2018} One TraDIS call per gene in each of
\goodallEnv{} media, \goodallGenes{} genes and \goodallLabels{} calls.
It is the simplest gene-level label in the host, and the one with three independent
calls to agree with, from the Keio construction outcome (Baba 2006), genetic
footprinting (Gerdes 2003) and the TnSeq no-insertion list (Price 2018). Genes are
keyed by \goodallIdentifier{}.

\pfstep{Chemical genomics, Nichols 2011} A colony-size S-score for every Keio deletion
in \nicholsEnv{} conditions: \nicholsGenes{} genes, \nicholsLabels{} scores. It is the
widest condition panel on one gene set in the host, so a held-out gene's whole profile
is the target. The RB-TnSeq compendium of Price 2018 poses the same task with a
sequencing readout on \priceEnv{} conditions and is the replication arm. Genes are keyed by
\nicholsIdentifier{}.

\pfstep{Proteome, Schmidt 2016} Absolute copies per cell for \schmidtGenes{} proteins
in \schmidtEnv{} growth conditions, \schmidtLabels{} values, on an unperturbed wild
type. The input is the coding sequence of the protein being measured, so this is the
row where a codon model is expected, untested here, to separate most from a one-hot
control. Mori 2021 measures the same quantity on a different platform and seven
samples; Caglar 2017 is the REL606 panel. Genes are keyed by \schmidtIdentifier{}.

\pfstep{Gene interaction, Babu 2014} A signed S score per deletion pair,
\babuLabels{} pairs over \babuGenes{} genes. It is the only pairwise label at scale in
the host and asks whether two gene vectors compose, which no single-gene task does.
Two splits are worth scoring: both pair members unseen, and one member unseen.
Butland 2008 is the unfiltered matrix of the same screens. Genes are keyed by
\babuIdentifier{}.

\pfstep{Second tier} Li 2014 gives an absolute protein synthesis rate per gene from
ribosome profiling in \liEnv{} media (\liGenes{} genes), the phenotype closest to what
a codon model reads. PRECISE-1K (Lamoureux 2023) gives a TPM per gene in \preciseEnv{}
distinct environments over \preciseInstances{} libraries (\preciseGenes{} genes), the
transcript-level counterpart of the proteome pick. Campos 2018 gives a vector of
cell-shape and cell-cycle features per Keio deletion (\camposGenes{} genes), the one
imaging phenotype in the host. Wang 2018 gives a gene-level CRISPRi fitness per gene
and screen (\wangGenes{} genes, \wangLabels{} values), the knockdown counterpart of
the deletion screens, and the one that reaches essential genes.

%% GENERATED FILE -- do not hand-edit.
%% SOURCE: experiments/database/scripts/build_ecoli_genslm_benchmark_table.py
%% off experiments/database/results/pre-build/2026-10-21/ecoli_datasets.json
\begin{landscape}
\begingroup
\scriptsize
\setlength{\tabcolsep}{3pt}
\begin{longtable}{@{}r >{\raggedright\arraybackslash}p{25mm} >{\raggedright\arraybackslash}p{20mm} >{\raggedright\arraybackslash}p{29mm} r r r >{\raggedright\arraybackslash}p{25mm} >{\raggedright\arraybackslash}p{23mm} >{\raggedright\arraybackslash}p{29mm} >{\raggedright\arraybackslash}p{46mm}@{}}
\caption{The short list for a gene-sequence encoding benchmark. \emph{Genes}, \emph{Labels}, \emph{Env} and \emph{Gene identifier} are read from a full scan of each store: distinct gene ids carrying a label, gene-attached numeric or boolean values, distinct stored environments, and the namespace the store keys genes by. \emph{Task}, \emph{Label}, \emph{Split}, \emph{Same task} and \emph{Why} are the pick list in the generating script.}\label{tab:shortlist}\\
\toprule
\# & Dataset & Task & Label & Genes & Labels & Env & Gene identifier & Split & Same task & Why \\
\midrule
\endfirsthead
\multicolumn{11}{@{}l}{\footnotesize\emph{Table~\ref{tab:shortlist}, continued}}\\
\toprule
\# & Dataset & Task & Label & Genes & Labels & Env & Gene identifier & Split & Same task & Why \\
\midrule
\endhead
\multicolumn{11}{@{}l}{\textbf{Core}} \\
1 & Gene\allowbreak{}Essentiality\allowbreak{}Goodall2018 & binary classification & essential or not, one TraDIS call per gene and medium & 4,181 & 8,309 & 2 & BW25113 locus tag, released as gene symbol & held-out genes & Gene\allowbreak{}Essentiality\allowbreak{}Baba2006, Gene\allowbreak{}Essentiality\allowbreak{}Gerdes2003, Gene\allowbreak{}Essentiality\allowbreak{}Price2018\allowbreak{}Ecoli & The simplest gene-level label in the host, and the one with three independent calls to agree with (Keio construction, genetic footprinting, TnSeq). A classifier on one vector per gene is the whole model. \\
2 & Env\allowbreak{}Chemgen\allowbreak{}Nichols2011 & regression, gene x condition & a colony-size S-score per (deletion, condition) & 3,949 & 1,286,962 & 324 & BW25113 locus tag, released as ECK accession & held-out genes, every condition & Rb\allowbreak{}Tnseq\allowbreak{}Price2018\allowbreak{}Ecoli, Env\allowbreak{}Chemgen\allowbreak{}Shiver2016, Carbon\allowbreak{}Source\allowbreak{}Tong2020, Env\allowbreak{}Chemgen\allowbreak{}Girgis2009 & The widest condition panel on one gene set. A held-out gene's whole profile is the target, and the condition enters as a second input the gene encoding does not have to carry. \\
3 & Proteome\allowbreak{}Schmidt2016 & regression, protein x condition & absolute copies per cell of one protein in one growth condition & 2,359 & 48,335 & 21 & gene symbol (readout key) & held-out genes, every condition & Proteome\allowbreak{}Mori2021, Proteome\allowbreak{}Caglar2017, Proteome\allowbreak{}Ishii2007 & A molecular level of the gene's own product, unperturbed, across conditions. Expected (untested) to be where a codon model separates from a one-hot control, because the input is the coding sequence whose product is measured. \\
4 & Gene\allowbreak{}Interaction\allowbreak{}Babu2014 & regression, gene pair & a signed colony-size S score per (query, array) deletion pair & 3,880 & 42,705 & 1 & MG1655 b-number, released as Keio JW number; b-number with a Keio sub-number; gene symbol & held-out genes (both pair members unseen), and one member unseen & Gene\allowbreak{}Interaction\allowbreak{}Butland2008, Crispri\allowbreak{}Cross\allowbreak{}Rachwalski2024 & The only pairwise label at scale in the host. It asks whether two gene vectors compose, which no single-gene task does. \\
\addlinespace
\multicolumn{11}{@{}l}{\textbf{Second tier}} \\
5 & Protein\allowbreak{}Synthesis\allowbreak{}Rate\allowbreak{}Li2014 & regression, gene x medium & absolute protein synthesis rate per gene from ribosome profiling & 3,410 & 8,644 & 3 & gene symbol (readout key) & held-out genes & Protein\allowbreak{}Turnover\allowbreak{}Gupta2024 & The phenotype closest to what a codon model reads: translation of the coding sequence itself, in three media. \\
6 & Rnaseq\allowbreak{}Lamoureux2023 & regression, gene x library & TPM of one gene in one RNA-seq library & 4,257 & 4,405,995 & 119 & b-number (readout key) & held-out genes, every library & Rnaseq\allowbreak{}Public\allowbreak{}K12\allowbreak{}Lamoureux2023, Rnaseq\allowbreak{}Caglar2017 & The transcript-level counterpart of the proteome pick, over a thousand libraries, so the per-gene target is a long expression profile. \\
7 & Morphology\allowbreak{}Campos2018 & multi-output regression & a vector of cell-shape and cell-cycle features per deletion & 3,664 & 53,508 & 1 & BW25113 locus tag & held-out genes & Growth\allowbreak{}Rate\allowbreak{}Campos2018 & The one imaging phenotype in the host: a different label family from growth and abundance, on the same Keio gene set. \\
8 & Crispri\allowbreak{}Gene\allowbreak{}Fitness\allowbreak{}Wang2018 & regression, gene x screen & gene-level CRISPRi fitness per (gene, screen) & 4,223 & 20,158 & 5 & MG1655 b-number & held-out genes & Crispri\allowbreak{}Knockdown\allowbreak{}Cui2018, Crispri\allowbreak{}Chemgen\allowbreak{}Choe2025, Mismatch\allowbreak{}Crispri\allowbreak{}Fitness\allowbreak{}Hawkins2020 & The knockdown counterpart of the deletion screens, already aggregated to the gene by its authors, so it covers essential genes the deletion collections cannot. \\
\bottomrule
\end{longtable}
\endgroup
\end{landscape}

